% Page 012

\seite{12}

der Kirche zu Seffern geschenkten Güter von der Els\ob
von derselben gemeinschaftlich zu Lehen hatten. Noch\ob
im Jahre 1811 erklärten beide Häuser in unserem\ob
Dörfchen, \enquote{Trauten Ann} und \enquote{Kotts Jannes},\ob
der Pfarrkirche die oben erwähnten 2 Malter Hafer\ob
Grüner Weizen jährlich zu liefern. Später\ob
unterließen beide Häuser die Lieferung, wahrschein\ohb
lich, weil sie einmal von französischer Cultur beeinflußt\ob
worden waren." Gerade sind beide Häuser Landbewohner\unsicher.

\pend
\abschnitt{Die Schule}
\pstart

Die Pfarrei Seffern hatte bis zum Jahre 1806 nur bloß\ob
zwei Kaplane, welche auch den Unterricht der Jugend\ob
\margmark{die Pfarrschüler}%
ertheilten. Gegen das 18.\ Jahrhundert erhielten\ob
diese Geistlichen meistens weltliche Gehülfen aus der\ob
Pfarrei. Das Schulhaus befand sich von jeher dicht\ob
an der Seite der Kirche auf der linken Pfade, wo\ob
auch die Kaplänei war, welche jetzt von den Ge\ohb
schwistern Henrion bewohnt wird. Die Pfarre zu\ob
Seffern begreifte die drei Filialen Sefferweich,\ob
\margmark{Die früheren\\Schülerverhältnisse}%
Schleid und Heilenbach. Während des Winters\ob
schickte der Pfarrer auf die Filialen begabte und\ob
fleißige, frühere Schüler, um dort zu schulmeistern,\ob
so gut sie es vermochten. In der zweiten Hälfte\ob
des 18.\ Jahrhunderts fand der Küster in der Pfarr-
\pend
